% =====================================================================
% CHƯƠNG 4 - PHÂN TÍCH VÀ THIẾT KẾ NỀN TẢNG EDGE AI
% =====================================================================
\chapter{PHÂN TÍCH VÀ THIẾT KẾ NỀN TẢNG EDGE AI}

Sau khi xác định pipeline tối ưu mô hình, chương này phân tích yêu cầu và thiết
kế nền tảng đưa model vào vận hành trên nhiều thiết bị biên. Nền tảng được tổ
chức theo mô hình Cloud Control Plane -- Edge Data Plane: Cloud giữ trạng thái
mong muốn (desired state) của toàn đội thiết bị, còn Edge Agent trên mỗi node
hội tụ trạng thái thực tế (actual state) về trạng thái mong muốn và báo cáo ngược
lại. Mã nguồn được chia thành các lớp domain, application và infrastructure theo
tinh thần Domain-Driven Design; hợp đồng dữ liệu giữa Cloud và Edge được đặt
trong một gói dùng chung để hai phía không bao giờ lệch schema.

\section{Yêu cầu hệ thống}

\subsection{Tác nhân và phạm vi}

Hệ thống có bốn nhóm tác nhân. \textbf{Người vận hành} thao tác qua Dashboard
hoặc trò chuyện với AI Assistant. \textbf{Cloud Control Plane} lưu trạng thái nền
tảng và phát lệnh. \textbf{Edge Agent} chạy thường trực trên mỗi Edge Node, nhận
lệnh và điều khiển container. \textbf{AI Runtime} là container thực thi mô hình
YOLO trên luồng camera và báo cáo chỉ số cho Edge Agent. Tài nguyên được tổ
chức phân cấp: một Cloud quản lý nhiều Edge Cluster; mỗi Cluster gồm nhiều Edge
Node; mỗi Node có đúng một Edge Agent; mỗi Agent quản lý nhiều AI Runtime.

\subsection{Yêu cầu chức năng}

\begin{longtable}{L{1.4cm}L{3.6cm}L{9.5cm}}
\caption{Yêu cầu chức năng của nền tảng}\label{tab:functional-requirements}\\
\toprule
\textbf{Mã} & \textbf{Nhóm} & \textbf{Mô tả}\\
\midrule\endfirsthead
\toprule
\textbf{Mã} & \textbf{Nhóm} & \textbf{Mô tả}\\
\midrule\endhead
FR01 & Quản lý Cluster & Tạo, liệt kê, xem chi tiết Edge Cluster; tên cluster là duy nhất.\\
FR02 & Quản lý Node & Đăng ký node vào cluster (tên duy nhất trong cluster), xem thông tin phần cứng, phiên bản Agent, trạng thái và thời điểm heartbeat cuối.\\
FR03 & Theo dõi liveness & Tự động chuyển node sang OFFLINE khi quá thời gian heartbeat (mặc định 30~s).\\
FR04 & Deployment & Triển khai một AI Runtime (image, model, cấu hình, giới hạn tài nguyên) lên một node; cập nhật từng trường; gỡ bỏ; theo dõi trạng thái PENDING $\rightarrow$ IN\_PROGRESS $\rightarrow$ SUCCEEDED/FAILED, REMOVING $\rightarrow$ REMOVED.\\
FR05 & Điều khiển Runtime & Start, stop, restart một runtime từ Cloud hoặc từ API cục bộ của Agent.\\
FR06 & Phân phối model & Tham chiếu model theo kho model (Hugging Face Hub) được ghim theo commit; Agent tải, lưu cache và gắn file model vào container.\\
FR07 & Telemetry & Thu thập chỉ số hệ thống, GPU và runtime (FPS, latency p50/p95/p99, hàng đợi, camera, confidence) theo chu kỳ; lưu mẫu mới nhất và lịch sử ngắn hạn.\\
FR08 & Event & Ghi nhận sự kiện runtime, lỗi và kết quả deployment theo từng node.\\
FR09 & Realtime & Đẩy trạng thái, telemetry và event tới Dashboard qua WebSocket theo bộ lọc topic.\\
FR10 & AI Assistant & Hỏi đáp bằng ngôn ngữ tự nhiên về cluster, node, runtime, deployment, telemetry, event; sinh biểu đồ và gợi ý câu hỏi tiếp theo; thao tác thay đổi chỉ thực hiện sau xác nhận.\\
FR11 & Giám sát hạ tầng & Thu metric node, container, GPU tại Edge và đẩy về Prometheus trên Cloud; cảnh báo theo luật.\\
\bottomrule
\end{longtable}

\subsection{Yêu cầu phi chức năng}

\begin{itemize}
  \item \textbf{Tính tự trị tại biên:} AI Runtime tiếp tục suy luận khi mất kết nối
        Cloud; Agent lưu desired state xuống đĩa và khôi phục sau khi khởi động lại.
  \item \textbf{Idempotency:} gửi lại cùng một deployment không tạo runtime mới;
        Agent bỏ qua lệnh trùng \texttt{command\_id}; kết quả của lệnh đã bị thay
        thế không được ghi đè trạng thái mới.
  \item \textbf{Tự phục hồi:} vòng reconcile định kỳ tạo lại container bị dừng bất
        thường, có giới hạn số lần thất bại liên tiếp để tránh vòng lặp khởi động.
  \item \textbf{Bảo mật:} REST, WebSocket và MCP yêu cầu API key; Agent cũng yêu
        cầu API key cho API cục bộ; bảng khóa chỉ lưu hàm băm; model gắn vào
        container ở chế độ chỉ đọc \cite{nist800207}.
  \item \textbf{Khả năng quan sát:} mọi dịch vụ ghi log có cấu trúc (tùy chọn
        JSON), cung cấp \texttt{/health} kiểm tra database, Redis và broker.
  \item \textbf{Khả năng mở rộng:} Cloud API không giữ trạng thái phiên; phần tiêu
        thụ MQTT có thể chạy trong cùng tiến trình hoặc tách thành worker riêng.
  \item \textbf{Tính di động:} Agent build đa kiến trúc (amd64, arm64) để chạy
        trên máy x86 và Jetson; runtime chọn nguồn khung hình và detector qua cấu hình.
\end{itemize}

\section{Kiến trúc tổng thể}

Kiến trúc tổng thể được thể hiện trong Hình~\ref{fig:platform-architecture}.
Người dùng làm việc qua Web Dashboard; mọi yêu cầu đi qua reverse proxy Nginx
tới Cloud API, WebSocket hoặc AI Assistant. Cloud Server là một ứng dụng FastAPI
duy nhất phục vụ đồng thời REST \texttt{/api/v1}, WebSocket \texttt{/api/v1/ws},
MCP server \texttt{/mcp} và bộ tiêu thụ MQTT (embedded worker). Trạng thái bền
vững nằm trong PostgreSQL, telemetry ngắn hạn nằm trong Redis, còn kênh Cloud--Edge
là MQTT broker. Model được lưu trên Hugging Face Hub; image runtime nằm trên
GitHub Container Registry (GHCR).

\begin{figure}[H]
\centering
\resizebox{\textwidth}{!}{%
\begin{tikzpicture}[font=\small,
  svc/.style={draw, rounded corners, align=center, minimum height=1cm, text width=2.8cm, fill=blue!5},
  store/.style={draw, align=center, minimum height=0.9cm, text width=2.4cm, fill=orange!8},
  ext/.style={draw, dashed, rounded corners, align=center, minimum height=0.9cm, text width=2.6cm, fill=gray!6},
  edge/.style={draw, rounded corners, align=center, minimum height=1.6cm, text width=3.6cm, fill=purple!5},
  lbl/.style={font=\scriptsize, fill=white, inner sep=1pt}]
  % Hàng người dùng và Cloud
  \node[svc, fill=green!6] (user) at (0,1.6) {Người vận hành\\(trình duyệt)};
  \node[svc] (web) at (0,-0.6) {Web Dashboard\\React + Nginx};
  \node[svc, text width=4.3cm, minimum height=2.9cm] (cloud) at (5.2,-0.6)
      {\textbf{Cloud Server} (FastAPI)\\REST /api/v1\\WebSocket /api/v1/ws\\MCP /mcp\\Embedded MQTT worker};
  \node[svc] (assistant) at (5.2,2.4) {AI Assistant\\LangGraph + LLM};
  % Kho dữ liệu
  \node[store] (pg) at (10.0,0.8) {PostgreSQL 16};
  \node[store] (redis) at (10.0,-0.6) {Redis 7};
  \node[store] (mqtt) at (10.0,-2.0) {MQTT Broker\\Mosquitto 2};
  % Dịch vụ ngoài
  \node[ext, text width=3.0cm] (ghcr) at (13.9,1.0) {GHCR (image runtime)};
  \node[ext, text width=3.0cm] (hf) at (13.9,-0.8) {Hugging Face Hub (kho model)};
  % Edge
  \node[edge] (edge1) at (6.0,-5.4) {\textbf{Edge Node 1}\\Edge Agent\\AI Runtime (container)\\Camera};
  \node[edge] (edge2) at (11.0,-5.4) {\textbf{Edge Node N}\\Edge Agent\\AI Runtime (container)\\Camera};

  \draw[flowarrow] (user) -- (web);
  \draw[flowarrow] (web) -- node[lbl, below]{HTTPS/WS} (cloud);
  \draw[flowarrow] (user.east) -- node[lbl, above]{SSE} (assistant.west);
  \draw[flowarrow] (assistant) -- node[lbl, right]{MCP} (cloud);
  \draw[flowarrow] (cloud.east |- pg) -- (pg);
  \draw[flowarrow] (cloud.east |- redis) -- (redis);
  \draw[flowarrow, <->] (cloud.east |- mqtt) -- (mqtt);
  \draw[flowarrow, <->] (mqtt.south) -- ++(0,-1.0) -| node[lbl, pos=0.25, above]{MQTT} (edge1.north);
  \draw[flowarrow, <->] (mqtt.south) -- ++(0,-1.0) -| (edge2.north);
  \draw[flowarrow, dashed] (hf.south) -- (13.9,-3.9) -- node[lbl, above]{tải model} (12.55,-3.9) -- (12.55,-4.6);
  \draw[flowarrow, dashed] (ghcr.east) -- (16.0,1.0) -- (16.0,-6.9) -- node[lbl, below, pos=0.5]{pull image runtime} (6.0,-6.9) -- (edge1.south);
\end{tikzpicture}}
\caption{Kiến trúc tổng thể nền tảng Edge AI}
\label{fig:platform-architecture}
\end{figure}

\subsection{Control Plane và Data Plane}

\textbf{Control Plane} (Cloud) chịu trách nhiệm quản lý tài nguyên, lưu desired
state của từng runtime dưới dạng Deployment, phát lệnh, tổng hợp trạng thái và
cung cấp dữ liệu giám sát. \textbf{Data Plane} (Edge) chịu trách nhiệm xử lý dữ
liệu tại chỗ: thu nhận khung hình, suy luận, quản lý vòng đời container và thu
thập telemetry. Hai mặt phẳng chỉ trao đổi qua các thông điệp đã định nghĩa trong
gói hợp đồng dùng chung \texttt{src/shared/contracts}.

Luồng điều khiển đi theo chiều \textbf{Client $\rightarrow$ Cloud $\rightarrow$
Edge Agent $\rightarrow$ AI Runtime}; luồng telemetry đi theo chiều ngược lại
\textbf{AI Runtime $\rightarrow$ Edge Agent $\rightarrow$ Cloud}. AI Runtime không
kết nối trực tiếp tới broker mà gửi chỉ số qua HTTP tới API cục bộ của Agent; nhờ
đó chỉ Agent giữ credential MQTT và Cloud chỉ cần tin cậy một thực thể trên mỗi node.

\subsection{Giao thức MQTT và cấu trúc topic}

MQTT \cite{mqtt5} được chọn làm kênh Cloud--Edge vì nhẹ, hỗ trợ kết nối từ sau
NAT (Edge chủ động kết nối ra ngoài), có cơ chế retained message và phù hợp với
đường truyền không ổn định. Topic được tổ chức theo phân cấp tài nguyên như
Bảng~\ref{tab:mqtt-topics}.

\begin{table}[H]
\centering
\caption{Cấu trúc topic MQTT giữa Cloud và Edge}
\label{tab:mqtt-topics}
\small
\begin{tabularx}{\textwidth}{L{6.6cm}L{2.5cm}Y}
\toprule
\textbf{Topic} & \textbf{Chiều} & \textbf{Nội dung}\\
\midrule
\path{edge/{cluster}/{node}/command/{type}} & Cloud$\rightarrow$Edge & Lệnh \texttt{deploy}, \texttt{update}, \texttt{start}, \texttt{stop}, \texttt{restart}, \texttt{remove}\\
\path{edge/{cluster}/{node}/state} & Edge$\rightarrow$Cloud & NodeState (retained), đồng thời là heartbeat\\
\path{edge/{cluster}/{node}/telemetry/{kind}} & Edge$\rightarrow$Cloud & \texttt{system}, \texttt{gpu}, \texttt{runtime}\\
\path{edge/{cluster}/{node}/event/{kind}} & Edge$\rightarrow$Cloud & \texttt{runtime}, \texttt{error}, \texttt{deployment}\\
\bottomrule
\end{tabularx}
\end{table}

Cloud đăng ký các wildcard \path{edge/+/+/state/#}, \path{edge/+/+/telemetry/+}
và \path{edge/+/+/event/+}; mỗi Agent chỉ đăng ký \path{edge/{cluster}/{node}/command/+}
của chính nó. Lớp \texttt{EdgeTopics} vừa sinh vừa phân tích topic, được cả Cloud
và Edge dùng chung.

\subsection{Tổ chức mã nguồn theo lớp}

Mã nguồn backend là một monorepo Python 3.12, tách phần ứng dụng khởi chạy
(\texttt{apps/}) khỏi phần nghiệp vụ (\texttt{src/}). Trong mỗi bounded context
(\texttt{cloud}, \texttt{edge}, \texttt{assistant}), mã được chia thành bốn lớp:

\begin{itemize}
  \item \textbf{domain}: entity, value object và interface repository, không phụ
        thuộc framework;
  \item \textbf{application}: service thực hiện use case, handler tiêu thụ thông
        điệp và \emph{port} (interface tới hạ tầng như container engine, camera,
        detector);
  \item \textbf{infrastructure}: \emph{adapter} hiện thực port và repository
        (PostgreSQL, Redis, Docker CLI, OpenCV, Ultralytics, MQTT);
  \item \textbf{container}: lớp lắp ráp phụ thuộc, chọn adapter theo cấu hình
        (ví dụ \texttt{storage\_backend=memory} để chạy không cần database).
\end{itemize}

Cách tổ chức port--adapter cho phép thay \texttt{DockerRuntimeDriver} bằng
\texttt{FakeRuntimeDriver}, hay \texttt{OpenCvFrameSource} bằng
\texttt{SyntheticFrameSource} khi kiểm thử, mà không sửa service nghiệp vụ.
Ba gói nội bộ dùng lại giữa các dịch vụ: \texttt{hopquangdo-messaging} (trừu
tượng PubSub trên MQTT), \texttt{hopquangdo-http} (xác thực API key, health
router, xử lý lỗi) và \texttt{hopquangdo-observability} (cấu hình log).

\section{Thiết kế Domain}

\subsection{Mô hình phân cấp tài nguyên}

Mô hình dữ liệu nghiệp vụ phản ánh phân cấp vật lý của hệ thống như
Hình~\ref{fig:resource-hierarchy}. Quan hệ Cluster--Node và Node--Runtime là
một--nhiều; Node--Agent là một--một nên Agent không cần entity riêng mà được
biểu diễn qua các thuộc tính của Node (\texttt{agent\_version},
\texttt{last\_heartbeat\_at}).

\begin{figure}[H]
\centering
\begin{tikzpicture}[node distance=0.45cm and 0.5cm]
  \node[groupbox] (cloud) {Cloud};
  \node[groupbox, below=of cloud] (cluster) {Edge Cluster};
  \node[groupbox, below=of cluster] (node) {Edge Node};
  \node[groupbox, below=of node] (agent) {Edge Agent};
  \node[groupbox, below=of agent] (runtime) {AI Runtime};
  \draw[flowarrow] (cloud) -- node[right, font=\scriptsize]{1 : N} (cluster);
  \draw[flowarrow] (cluster) -- node[right, font=\scriptsize]{1 : N} (node);
  \draw[flowarrow] (node) -- node[right, font=\scriptsize]{1 : 1} (agent);
  \draw[flowarrow] (agent) -- node[right, font=\scriptsize]{1 : N} (runtime);
\end{tikzpicture}
\caption{Phân cấp tài nguyên của nền tảng}
\label{fig:resource-hierarchy}
\end{figure}

\subsection{Cloud domain}

Bảng~\ref{tab:cloud-entities} tóm tắt các entity của Cloud. Entity được hiện thực
bằng \texttt{dataclass}, không phụ thuộc ORM; repository chuyển đổi giữa entity và
bản ghi database.

\begin{table}[H]
\centering
\caption{Các entity chính của Cloud domain}
\label{tab:cloud-entities}
\begin{tabularx}{\textwidth}{L{2.4cm}L{5.2cm}Y}
\toprule
\textbf{Entity} & \textbf{Thuộc tính chính} & \textbf{Quy tắc nghiệp vụ}\\
\midrule
Cluster & id, name, status, location, metadata & Tên duy nhất trên toàn hệ thống.\\
Node & id, cluster\_id, name, hostname, ip, status, agent\_version, cpu\_cores, memory\_total\_mb, gpu\_name, last\_heartbeat\_at & Tên duy nhất trong cluster; thông tin phần cứng do Agent tự báo cáo; OFFLINE khi quá hạn heartbeat.\\
Runtime & id, node\_id, name, type, status, image, model\_id, config, environment, resource\_limits, container\_id & Trạng thái do Edge báo cáo; bị xóa khi Edge xác nhận gỡ bỏ.\\
Deployment & id, cluster\_id, node\_id, runtime\_id, spec, status, last\_command\_id, message & Desired state của một runtime trên một node; chỉ chấp nhận kết quả của lệnh mới nhất.\\
\bottomrule
\end{tabularx}
\end{table}

\textbf{Deployment là aggregate trung tâm.} Mỗi lần phát lệnh, phương thức
\texttt{dispatched(command\_id)} ghi lại \texttt{last\_command\_id} và chuyển trạng
thái. Khi Edge báo kết quả, \texttt{apply\_result(command\_id, status)} chỉ áp dụng
nếu \texttt{command\_id} trùng lệnh mới nhất; kết quả của lệnh cũ đến muộn bị bỏ
qua. Đây là cơ chế chống ghi đè ngược trạng thái khi mạng giao lại thông điệp
không theo thứ tự. Deployment ở trạng thái REMOVED là trạng thái kết thúc và
không nhận lệnh mới.

\begin{figure}[H]
\centering
\resizebox{\textwidth}{!}{%
\begin{tikzpicture}[font=\small,
  st/.style={draw, rounded corners, align=center, minimum height=0.8cm, minimum width=2.9cm, fill=blue!5},
  lbl/.style={font=\scriptsize, fill=white, inner sep=1.5pt}]
  \node[st] (pending) at (0,0) {PENDING};
  \node[st] (prog) at (4.2,0) {IN\_PROGRESS};
  \node[st] (ok) at (8.6,1.4) {SUCCEEDED};
  \node[st] (fail) at (8.6,-1.4) {FAILED};
  \node[st] (removing) at (13.0,0) {REMOVING};
  \node[st, fill=gray!15] (removed) at (17.2,0) {REMOVED};
  \draw[flowarrow] (pending) -- node[lbl, above=2pt]{dispatch} (prog);
  \draw[flowarrow] (prog.north east) -- node[lbl, above left]{kết quả OK} (ok.west);
  \draw[flowarrow] (prog.south east) -- node[lbl, below left]{kết quả lỗi} (fail.west);
  \draw[flowarrow] (ok.east) -- node[lbl, above right]{remove} (removing.north west);
  \draw[flowarrow] (fail.east) -- node[lbl, below right]{remove} (removing.south west);
  \draw[flowarrow] (removing) -- node[lbl, above=2pt]{Edge xác nhận} (removed);
  \draw[flowarrow, dashed] (ok.north) -- ++(0,0.7) -| node[lbl, pos=0.25, above]{update / start / stop / restart} (prog.north);
\end{tikzpicture}}
\caption{Máy trạng thái của Deployment}
\label{fig:deployment-state-machine}
\end{figure}

Value object gồm \texttt{ModelRef} (repo\_id, filename, revision),
\texttt{ModelLocation} (URL tải đã ghim theo commit SHA) và
\texttt{NodeHeartbeat} (các thông tin Node trích từ bản tin NodeState).

\subsection{Edge domain}

Phía Edge có entity \texttt{RuntimeInstance} gồm \texttt{RuntimeSpec} mong muốn,
\texttt{desired\_running}, trạng thái hiện tại, \texttt{container\_id},
\texttt{restart\_count}, số lần khởi chạy thất bại liên tiếp và thông điệp lỗi gần
nhất. \texttt{DesiredStateRepository} lưu danh sách RuntimeInstance xuống tệp JSON
trong \texttt{state\_dir} để Agent khôi phục sau khi khởi động lại.

Vòng đời runtime dùng chung một enum giữa hai phía (Bảng~\ref{tab:runtime-states}).

\begin{table}[H]
\centering
\caption{Trạng thái của AI Runtime}
\label{tab:runtime-states}
\begin{tabularx}{\textwidth}{L{3cm}Y}
\toprule
\textbf{Trạng thái} & \textbf{Ý nghĩa}\\
\midrule
CREATING & Agent đang tải model và pull image.\\
STARTING & Container đang được khởi chạy.\\
RUNNING & Container chạy bình thường.\\
DEGRADED & Container chạy nhưng hiệu năng suy giảm.\\
STOPPING / STOPPED & Đang dừng / đã dừng theo yêu cầu.\\
RESTARTING & Đang khởi động lại.\\
FAILED & Khởi chạy thất bại hoặc vượt số lần thử tối đa.\\
UNKNOWN & Chưa xác định (ví dụ ngay sau khi Agent khởi động lại).\\
\bottomrule
\end{tabularx}
\end{table}

\subsection{Hợp đồng dùng chung Cloud--Edge}

Gói \texttt{src/shared/contracts} định nghĩa bằng Pydantic toàn bộ thông điệp
trao đổi: \texttt{Command} và \texttt{RuntimeSpec} (Cloud$\rightarrow$Edge);
\texttt{NodeState}, \texttt{RuntimeState}, \texttt{SystemTelemetry},
\texttt{GpuTelemetry}, \texttt{RuntimeTelemetry}, \texttt{EdgeEvent},
\texttt{DeploymentEvent} (Edge$\rightarrow$Cloud) cùng các enum trạng thái. Giá
trị enum trùng khớp với kiểu enum trong PostgreSQL. Vì hai phía import cùng một
mã nguồn, thông điệp sai schema bị từ chối ngay tại nơi nhận thay vì làm hỏng dữ
liệu.

\texttt{RuntimeSpec} là desired state của một runtime:
\texttt{runtime\_id}, \texttt{name}, \texttt{type} (VISION, SAFETY, ANOMALY, OCR,
CUSTOM), \texttt{image\_uri}, \texttt{image\_tag}, \texttt{model\_id},
\texttt{model\_uri}, \texttt{config}, \texttt{environment} và
\texttt{resource\_limits}. Deploy và update là idempotent theo \texttt{runtime\_id}.

\section{Thiết kế Cloud Control Plane}
\label{sec:cloud-design}

Cloud Control Plane gồm các application service xử lý yêu cầu đồng bộ từ người
dùng và các handler xử lý thông điệp bất đồng bộ từ Edge
(Hình~\ref{fig:cloud-components}).

\begin{figure}[H]
\centering
\resizebox{\textwidth}{!}{%
\begin{tikzpicture}[font=\small,
  b/.style={draw, rounded corners, align=center, minimum height=0.9cm, text width=3.3cm, fill=blue!5},
  h/.style={draw, rounded corners, align=center, minimum height=1.0cm, text width=3.3cm, fill=orange!8},
  r/.style={draw, align=center, minimum height=0.9cm, fill=gray!8},
  lbl/.style={font=\scriptsize, fill=white, inner sep=1.5pt}]
  % Tầng vào
  \node[b, fill=green!6, text width=16cm] (entry) at (6,3.2) {REST Router \qquad | \qquad WebSocket Router \qquad | \qquad MCP Tools};
  % Application service
  \node[b] (fleet) at (0,1.6) {FleetService};
  \node[b] (deploy) at (4,1.6) {DeploymentService};
  \node[b] (monitor) at (8,1.6) {MonitoringService};
  \node[b] (realtime) at (12,1.6) {RealtimeService};
  % MQTT
  \node[r, fill=purple!6, text width=15.6cm] (mqtt) at (6,-0.2) {MQTT Broker};
  % Handler
  \node[h] (state) at (0,-2.0) {NodeStateHandler};
  \node[h] (devt) at (4,-2.0) {DeploymentEvent\\Handler};
  \node[h] (tele) at (8,-2.0) {TelemetryHandler};
  \node[h] (evt) at (12,-2.0) {EdgeEventHandler};
  \node[h, fill=red!6, text width=3.7cm] (live) at (-4.2,-2.0) {NodeLiveness\\Service (quét 10 s)};
  % Lưu trữ
  \node[r, text width=10cm] (pg) at (2,-4.0) {PostgreSQL: clusters, nodes, runtimes, deployments, node\_logs};
  \node[r, text width=3.3cm] (redis) at (10,-4.0) {Redis: telemetry};

  \foreach \n in {fleet, deploy, monitor, realtime} \draw[flowarrow] (entry.south -| \n) -- (\n);
  \draw[flowarrow] (deploy) -- node[lbl, right]{CommandGateway: publish lệnh} (deploy |- mqtt.north);
  \draw[flowarrow, <-] (realtime) -- node[lbl, right]{edge/\#} (realtime |- mqtt.north);
  \foreach \n in {state, devt, tele, evt} \draw[flowarrow] (\n |- mqtt.south) -- (\n);
  \draw[flowarrow, dashed] (state) -- (state |- pg.north);
  \draw[flowarrow, dashed] (devt) -- (devt |- pg.north);
  \draw[flowarrow, dashed] (evt.south) -- ++(0,-0.5) -| ([xshift=2.5cm]pg.north);
  \draw[flowarrow, dashed] (tele) -- (tele |- redis.north);
  \draw[flowarrow, dashed] (live) |- (pg.west);
  \draw[flowarrow, dashed] (fleet.west) -| ([xshift=-0.3cm]live.north);
\end{tikzpicture}}
\caption{Các thành phần của Cloud Control Plane}
\label{fig:cloud-components}
\end{figure}

\subsection{Fleet Service}

FleetService hiện thực các use case trên phân cấp Cluster $\rightarrow$ Node
$\rightarrow$ Runtime: tạo cluster (từ chối tên trùng), đăng ký node vào cluster
(từ chối tên trùng trong cùng cluster), liệt kê và tra cứu. Node được đăng ký
trước trên Cloud; khi Agent với \texttt{node\_id} tương ứng bắt đầu gửi
NodeState, Cloud cập nhật thông tin phần cứng. Bản tin NodeState từ node chưa
đăng ký bị ghi log và bỏ qua, tránh việc thiết bị lạ tự chen vào hệ thống.

\subsection{Deployment Service}

DeploymentService ghi desired state rồi phát lệnh; kết quả được Edge báo về bất
đồng bộ. Do đó các API tạo, cập nhật, gỡ bỏ trả mã HTTP \texttt{202 Accepted}.
Các quy tắc chính:

\begin{itemize}
  \item \textbf{Tạo mới:} kiểm tra node tồn tại. Nếu đã có runtime cùng tên trên
        node với deployment còn hiệu lực: spec giống hệt thì trả lại deployment cũ
        (gửi lại idempotent); spec khác thì báo xung đột và yêu cầu dùng PUT.
        Ngược lại, sinh \texttt{runtime\_id}, tạo bản ghi Runtime (CREATING) và
        Deployment, rồi phát lệnh \texttt{deploy}.
  \item \textbf{Cập nhật:} chỉ ghi đè những trường được gửi lên, đồng bộ bản ghi
        Runtime và phát lệnh \texttt{update}.
  \item \textbf{Gỡ bỏ:} chuyển REMOVING và phát lệnh \texttt{remove}; bản ghi
        Deployment được giữ làm lịch sử, bản ghi Runtime bị xóa khi Edge xác nhận.
  \item \textbf{Điều khiển:} start/stop/restart chỉ áp dụng cho runtime có
        deployment còn hiệu lực.
\end{itemize}

Trước khi publish, service luôn gọi \texttt{dispatched()} và lưu deployment, nên
trạng thái IN\_PROGRESS và \texttt{last\_command\_id} đã bền vững trước khi lệnh
rời khỏi Cloud. CommandGateway đóng gói việc publish lên topic
\path{edge/{cluster}/{node}/command/{type}}.

\subsection{Các handler bất đồng bộ}

\begin{itemize}
  \item \textbf{NodeStateHandler} tiêu thụ \path{edge/+/+/state}: cập nhật
        heartbeat, trạng thái Agent, thông tin phần cứng và đồng bộ trạng thái,
        \texttt{container\_id} của từng runtime do node báo cáo.
  \item \textbf{DeploymentEventHandler} tiêu thụ \path{event/deployment}: tìm
        deployment theo \texttt{deployment\_id} hoặc \texttt{runtime\_id}, áp dụng
        kết quả qua \texttt{apply\_result}, cập nhật trạng thái runtime hoặc xóa
        runtime khi đã REMOVED.
  \item \textbf{TelemetryHandler} tiêu thụ \path{telemetry/{kind}}: kiểm tra schema
        theo loại, lưu mẫu vào Redis; mẫu runtime được khóa theo
        \texttt{runtime:\{runtime\_id\}}.
  \item \textbf{EdgeEventHandler} tiêu thụ \path{event/+}: ghi mọi sự kiện vào nhật
        ký của node (\texttt{node\_logs}).
  \item \textbf{NodeLivenessService} quét định kỳ (10~s) và chuyển các node có
        heartbeat cũ hơn \texttt{node\_heartbeat\_timeout\_sec} (30~s) sang OFFLINE.
\end{itemize}

\subsection{Monitoring và Realtime Service}

MonitoringService cung cấp mẫu telemetry mới nhất của một node (theo từng loại),
lịch sử gần đây của một loại và danh sách event. RedisTelemetryRepository lưu
\texttt{telemetry:\{node\}:latest} dạng hash và \texttt{telemetry:\{node\}:\{kind\}}
dạng list giới hạn 720 mẫu (tương đương 1 giờ với chu kỳ 5~s), có TTL 1 giờ.

RealtimeService đăng ký \path{edge/#} trên broker và phân phối thông điệp tới các
WebSocket client theo bộ lọc topic của từng client (hỗ trợ wildcard MQTT
\texttt{+} và \texttt{\#}). Client có thể đổi bộ lọc khi đang kết nối bằng thông
điệp \texttt{subscribe}/\texttt{unsubscribe}; xác nhận \texttt{ack} đi qua cùng
hàng đợi để giữ thứ tự với dữ liệu.

\subsection{Triển khai tiến trình}

Cloud chạy như một tiến trình duy nhất trên cổng 8000, phục vụ REST, WebSocket,
MCP và embedded worker (\texttt{CLOUD\_EMBEDDED\_WORKER=true}). Khi cần mở rộng,
worker có thể tách thành tiến trình riêng (\texttt{apps/cloud/worker}) trong khi API
được nhân bản. Chế độ \texttt{CLOUD\_STORAGE\_BACKEND=memory} thay PostgreSQL và
Redis bằng repository trong bộ nhớ, phục vụ phát triển và kiểm thử tự động.

\section{Thiết kế Edge Agent}

Edge Agent là tiến trình thường trực duy nhất được cài tĩnh trên Edge Node; mọi
AI Runtime do Agent tạo động từ image runtime. Các thành phần chính được thể hiện
trong Hình~\ref{fig:edge-agent-components}.

\begin{figure}[H]
\centering
\resizebox{0.95\textwidth}{!}{%
\begin{tikzpicture}[font=\small,
  b/.style={draw, rounded corners, align=center, minimum height=1.0cm, text width=3.9cm, fill=blue!5}]
  \node[b, fill=purple!6] (mqtt) at (0,4.4) {MQTT\\(command/+)};
  \node[b, fill=green!6] (api) at (4.4,4.4) {Local API :8081\\(REST)};
  \node[b] (cmd) at (2.2,2.6) {CommandHandler\\(chống lệnh trùng)};
  \node[b] (rt) at (2.2,0.6) {RuntimeService\\(apply / reconcile)};
  \node[b] (desired) at (-3.2,0.6) {JsonDesiredState\\Repository};
  \node[b] (drv) at (-0.2,-1.6) {DockerRuntimeDriver\\(Docker CLI)};
  \node[b] (fetch) at (4.4,-1.6) {HttpArtifactFetcher\\(cache model)};
  \node[b] (agent) at (8.2,0.6) {AgentService\\(vòng lặp định kỳ)};
  \node[b] (state) at (10.0,2.6) {StateReportService\\(10 s)};
  \node[b] (tele) at (10.0,-1.6) {TelemetryReportService\\(5 s)};
  \draw[flowarrow] (mqtt) -- (cmd);
  \draw[flowarrow] (api) -- (cmd);
  \draw[flowarrow] (cmd) -- (rt);
  \draw[flowarrow] (rt) -- (drv);
  \draw[flowarrow] (rt) -- (fetch);
  \draw[flowarrow, <->] (rt) -- (desired);
  \draw[flowarrow] (agent) -- node[font=\scriptsize, above=3pt]{reconcile 15 s} (rt);
  \draw[flowarrow] (agent) -- (state);
  \draw[flowarrow] (agent) -- (tele);
\end{tikzpicture}}
\caption{Các thành phần chính của Edge Agent}
\label{fig:edge-agent-components}
\end{figure}

\subsection{Vòng lặp điều khiển}

Khi khởi động, AgentService kết nối broker, khôi phục desired state từ đĩa,
đăng ký topic lệnh và phát ngay một NodeState. Sau đó Agent chạy ba vòng lặp độc
lập (chu kỳ cấu hình được):

\begin{itemize}
  \item \textbf{State} (10~s): phát NodeState dạng retained, đóng vai trò heartbeat
        và báo cáo toàn bộ runtime. Agent ở trạng thái DEGRADED nếu có runtime
        mong muốn chạy nhưng đang FAILED hoặc DEGRADED.
  \item \textbf{Telemetry} (5~s): thu và phát telemetry \texttt{system} và \texttt{gpu}.
  \item \textbf{Reconcile} (15~s): so sánh desired state với trạng thái container
        thực tế và hội tụ.
\end{itemize}

Khi dừng, Agent hủy các vòng lặp và phát NodeState OFFLINE trước khi ngắt kết nối.

\subsection{Xử lý lệnh}

CommandHandler nhận lệnh từ MQTT hoặc từ API cục bộ; hai nguồn đi qua cùng một
đường xử lý nên kết quả đều được báo về Cloud. Lệnh không đúng schema bị từ chối;
lệnh có \texttt{command\_id} đã xử lý (bộ nhớ đệm 256 id gần nhất) bị bỏ qua để
chống giao lặp của broker. Sau khi thực thi, Agent phát \texttt{DeploymentEvent}
chứa \texttt{command\_id}, \texttt{deployment\_id}, trạng thái deployment và trạng
thái runtime, rồi phát ngay một NodeState mới. Lỗi nghiệp vụ (\texttt{AppException})
và lỗi bất ngờ đều được chuyển thành kết quả FAILED có thông điệp, không làm dừng Agent.

\subsection{Desired state và reconcile}

RuntimeService giữ desired state của mọi runtime trên node và là nơi duy nhất
thay đổi container. Mọi thao tác đi qua một khóa bất đồng bộ nên các lệnh trên
cùng node được tuần tự hóa. Các thao tác đều idempotent:

\begin{itemize}
  \item \texttt{apply(spec)}: nếu spec giống hệt và runtime đang RUNNING thì không
        làm gì; nếu khác, hủy container cũ, lưu desired state và khởi chạy mới.
  \item \texttt{start}/\texttt{stop}: đặt \texttt{desired\_running} rồi thao tác
        container; \texttt{remove} với runtime không tồn tại vẫn thành công.
  \item \texttt{reconcile()}: với runtime mong muốn dừng mà container đang chạy thì
        dừng; mong muốn chạy mà container mất hoặc đã thoát thì ghi lý do (mã thoát,
        lỗi), xóa container cũ và tạo lại; sau \texttt{max\_failures} (5) lần thất bại
        liên tiếp thì đánh dấu FAILED và ngừng thử.
  \item \texttt{restore()}: sau khi Agent khởi động lại, đọc desired state từ đĩa và
        tìm container còn sống theo nhãn \texttt{edgeops.runtime\_id} để tiếp quản
        thay vì tạo trùng.
\end{itemize}

Mỗi lần trạng thái runtime thay đổi, Agent phát một \texttt{EdgeEvent}
(\texttt{runtime.status}); chuyển sang FAILED được phát trên kênh \texttt{error}
với mức ERROR.

\subsection{Khởi chạy runtime}

Quy trình \texttt{\_launch} gồm: (1) chuyển CREATING; (2) nếu spec có
\texttt{model\_uri}, tải model qua ArtifactFetcher và gắn vào container tại
\path{/models/<file>} ở chế độ chỉ đọc; (3) pull image theo chính sách
\texttt{always}/\texttt{missing}/\texttt{never}; (4) chuyển STARTING và chạy
container với biến môi trường \texttt{RUNTIME\_ID}, \texttt{RUNTIME\_NAME},
\texttt{RUNTIME\_CONFIG}, \texttt{MODEL\_PATH}, \texttt{NODE\_ID},
\texttt{AGENT\_URL}; (5) chuyển RUNNING. Giới hạn tài nguyên trong spec được ánh
xạ sang tham số Docker: \texttt{--gpus}, \texttt{--runtime nvidia} (Jetson),
\texttt{--cpus}, \texttt{--memory}, \texttt{--device} (camera). Container được gắn
chính sách \texttt{--restart unless-stopped} nên Docker tự khởi động lại khi
thiết bị reboot.

\subsection{API cục bộ}

Agent cung cấp REST API trên cổng 8081 (bảo vệ bằng API key) gồm hai nhóm:
nhóm quản lý (\texttt{GET /status}, \texttt{GET/POST/PUT/DELETE /runtimes},
\texttt{POST /runtimes/\{id\}/\{start|stop|restart\}}) cho phép vận hành tại chỗ
khi không có Cloud; và nhóm báo cáo (\texttt{POST /runtimes/\{id\}/telemetry},
\texttt{POST /runtimes/\{id\}/events}) để AI Runtime gửi chỉ số. Agent làm giàu mẫu
telemetry runtime với \texttt{node\_id}, trạng thái, số lần restart, tính
\texttt{effective\_fps\_ratio} = FPS thực/FPS mục tiêu, rồi chuyển tiếp lên Cloud.

\section{Thiết kế Edge Runtime}

AI Runtime là một container độc lập do Agent tạo ra. Toàn bộ cấu hình được
truyền qua biến môi trường; khóa \texttt{RUNTIME\_CONFIG} (JSON) chọn nguồn
khung hình (\texttt{source}), FPS mục tiêu (\texttt{target\_fps}, mặc định 15),
detector (\texttt{ultralytics} khi có \texttt{MODEL\_PATH}, ngược lại
\texttt{synthetic}), ngưỡng tin cậy (\texttt{conf}) và chu kỳ báo cáo. Luồng xử
lý được thể hiện trong Hình~\ref{fig:edge-runtime-design}.

\begin{figure}[H]
\centering
\resizebox{\textwidth}{!}{%
\begin{tikzpicture}[node distance=0.55cm and 0.7cm]
  \node[flowbox] (camera) {Camera\\RTSP/USB/file};
  \node[flowbox, right=of camera] (capture) {Capture task\\(FrameSource)};
  \node[flowbox, right=of capture, fill=orange!8] (queue) {Hàng đợi\\8 khung};
  \node[flowbox, right=of queue] (infer) {Inference task\\(Detector)};
  \node[flowbox, right=of infer, text width=3cm] (metrics) {Inference\\Metrics};
  \node[flowbox, below=1.4cm of metrics, text width=3cm] (report) {Report task\\(5 s)};
  \node[flowbox, left=of report, fill=green!6] (agent) {Edge Agent\\HTTP API};
  \draw[flowarrow] (camera) -- (capture);
  \draw[flowarrow] (capture) -- (queue);
  \draw[flowarrow] (queue) -- (infer);
  \draw[flowarrow] (infer) -- (metrics);
  \draw[flowarrow] (metrics) -- (report);
  \draw[flowarrow] (report) -- node[font=\scriptsize, above]{telemetry} (agent);
  \draw[flowarrow, dashed] (capture.south) -- ++(0,-0.6) -| node[pos=0.25, above, font=\scriptsize]{frame drop, sự kiện camera} ([xshift=-0.6cm]metrics.south);
\end{tikzpicture}}
\caption{Thiết kế pipeline Edge Runtime}
\label{fig:edge-runtime-design}
\end{figure}

\textbf{Ba tác vụ song song.} InferenceService chạy ba tác vụ trong một
\texttt{TaskGroup}: \emph{capture} đọc khung từ FrameSource; \emph{infer} lấy khung
từ hàng đợi và gọi Detector với thời gian chờ tối đa 2~s; \emph{report} định kỳ
chụp ảnh chỉ số và gửi cho Agent. Khi một tác vụ lỗi nghiêm trọng, cả nhóm dừng
và container thoát, để vòng reconcile của Agent tạo lại container.

\textbf{Hàng đợi giới hạn, ưu tiên khung mới.} Hàng đợi chỉ giữ 8 khung. Khi
đầy, khung cũ nhất bị loại và được tính là \emph{frame drop}. Chính sách này giữ
độ trễ không tăng vô hạn khi detector chậm hơn camera: hệ thống chấp nhận bỏ
khung để luôn suy luận trên dữ liệu mới nhất.

\textbf{Nguồn khung và detector là port.} \texttt{OpenCvFrameSource} đọc RTSP,
USB (chỉ số thiết bị) hoặc tệp video, tự kết nối lại sau 3~s khi luồng mất và
phát sự kiện \texttt{camera.disconnected}. \texttt{UltralyticsDetector} nạp mô
hình qua API \texttt{YOLO(...)} của Ultralytics nên dùng được trực tiếp tệp
\texttt{.pt}, \texttt{.onnx} hoặc TensorRT \texttt{.engine}; lời gọi suy luận
chạy trong thread riêng để không chặn vòng sự kiện. Các adapter
\texttt{SyntheticFrameSource} và \texttt{SyntheticDetector} cho phép chạy toàn bộ
pipeline khi không có camera hoặc GPU.

\textbf{Chỉ số runtime.} \texttt{InferenceMetrics} tích lũy trong mỗi cửa sổ báo
cáo: FPS suy luận và FPS camera, latency trung bình và p50/p95/p99, tỷ lệ lỗi và
timeout, tỷ lệ frame drop, tuổi khung mới nhất, thời gian camera bị đứng hình,
độ sáng trung bình, số detection, confidence trung bình/độ lệch/nhỏ nhất và
entropy phân bố lớp. Các trường này khớp tên với đặc tả telemetry ở
Mục~\ref{sec:monitoring-design}.

\section{Thiết kế Model Deployment}

\subsection{Lưu trữ và phân phối model}

Model sau tối ưu (Chương~\ref{chap:toi-uu}) được lưu trên một kho model kiểu
Hugging Face Hub, có thể là huggingface.co, mirror hoặc máy chủ tự dựng. Mỗi
model được tham chiếu bằng \texttt{ModelRef(repo\_id, filename, revision)}.
Khi phân giải, \texttt{HuggingFaceModelRepository} chuyển \texttt{revision}
(nhánh hoặc tag) thành \emph{commit SHA} cụ thể và sinh URL tải đã ghim theo SHA
đó (\texttt{ModelLocation}). Vì URL trỏ đến nội dung bất biến của một commit, hai
node triển khai cùng một deployment chắc chắn nhận cùng một tệp model, kể cả khi
nhánh \texttt{main} được cập nhật giữa chừng.

Phía Edge, \texttt{HttpArtifactFetcher} tải model qua HTTPS theo từng khối 1~MB
vào tệp tạm \texttt{.part}, sau đó đổi tên nguyên tử thành tệp đích. Tệp được lưu
cache theo hàm băm SHA-256 của URL (bỏ phần query), nên triển khai lại cùng model
không tải lại. URI \texttt{file://} hoặc đường dẫn cục bộ được dùng trực tiếp,
thuận tiện cho kiểm thử trên thiết bị.

\subsection{Luồng triển khai}

Luồng triển khai đầy đủ được mô tả trong Hình~\ref{fig:model-deployment-design}.

\begin{figure}[H]
\centering
\resizebox{\textwidth}{!}{%
\begin{tikzpicture}[font=\small,
  msg/.style={font=\scriptsize, above, fill=white, inner sep=1pt},
  note/.style={font=\scriptsize, align=left, anchor=west, draw=gray, fill=yellow!8, rounded corners=2pt, inner sep=2pt}]
  \foreach \x/\n/\t in {0/u/{Dashboard hoặc Assistant}, 3.4/c/{Cloud API}, 6.8/m/{MQTT}, 10.2/a/{Edge Agent}, 13.6/r/{Docker, AI Runtime}}
    { \node[draw, rounded corners, fill=blue!5, align=center, text width=2.6cm, minimum height=1.1cm] (\n) at (\x,0) {\t};
      \draw[dashed, gray] (\x,-0.6) -- (\x,-11.2); }
  \draw[flowarrow] (0,-1.3) -- node[msg]{POST /deployments} (3.4,-1.3);
  \node[note] at (3.5,-2.0) {lưu Runtime (CREATING),\\Deployment (IN\_PROGRESS)};
  \draw[flowarrow] (3.4,-2.9) -- node[msg]{command/deploy} (6.8,-2.9);
  \draw[flowarrow, dashed] (3.4,-3.5) -- node[msg]{202 Accepted} (0,-3.5);
  \draw[flowarrow] (6.8,-4.1) -- node[msg]{command/deploy} (10.2,-4.1);
  \node[note] at (10.3,-4.9) {chống trùng, lưu desired state,\\tải model (có cache)};
  \draw[flowarrow] (10.2,-5.9) -- node[msg]{pull image, run} (13.6,-5.9);
  \draw[flowarrow, dashed] (13.6,-6.6) -- node[msg]{container id} (10.2,-6.6);
  \draw[flowarrow] (10.2,-7.3) -- node[msg]{event/deployment} (6.8,-7.3);
  \draw[flowarrow] (6.8,-7.9) -- node[msg]{DeploymentEvent} (3.4,-7.9);
  \node[note] at (3.5,-8.7) {apply\_result: SUCCEEDED,\\Runtime RUNNING};
  \draw[flowarrow] (10.2,-9.6) -- node[msg]{state (retained)} (6.8,-9.6);
  \draw[flowarrow] (6.8,-10.2) -- node[msg]{NodeState} (3.4,-10.2);
  \draw[flowarrow, dashed] (3.4,-10.8) -- node[msg]{WebSocket} (0,-10.8);
\end{tikzpicture}}
\caption{Trình tự triển khai model từ Dashboard đến AI Runtime}
\label{fig:model-deployment-design}
\end{figure}

\subsection{Cập nhật và khôi phục}

Cập nhật model hoặc cấu hình là một lệnh \texttt{update} với spec mới trên cùng
\texttt{runtime\_id}. Agent chỉ tạo lại container khi spec thực sự thay đổi. Nếu
khởi chạy với model mới thất bại, Agent báo FAILED kèm thông điệp lỗi (ví dụ
không tải được model, image không tồn tại); người vận hành hoặc Assistant có thể
gửi lại một \texttt{update} trỏ về model trước đó để khôi phục. Vì mỗi lệnh
update mang \texttt{command\_id} mới, kết quả của lần thử lỗi trước không ghi đè
kết quả khôi phục. Việc tự động canary, chia đợt và rollback theo ngưỡng telemetry
được xác định là hướng phát triển (Chương~\ref{chap:ket-luan}).

\section{Thiết kế Monitoring và cảnh báo}
\label{sec:monitoring-design}

Giám sát được thiết kế theo hai kênh bổ trợ nhau. Kênh \textbf{telemetry ứng
dụng} đi qua MQTT, mang ngữ nghĩa của nền tảng (runtime, FPS, camera, chất lượng
đầu ra model) và phục vụ Dashboard, Assistant. Kênh \textbf{metric hạ tầng} dùng
hệ sinh thái Prometheus \cite{prometheus} để thu metric node, container và GPU một cách chuẩn hóa,
lưu dài hạn và cảnh báo theo luật.

\subsection{Đặc tả telemetry}

Đặc tả telemetry chia chỉ số thành bảy nhóm và ba loại thông điệp
(Bảng~\ref{tab:telemetry-groups}). Tất cả trường đều tùy chọn để thiết bị chỉ báo
cáo những gì đo được; tên trường thống nhất giữa Edge Agent, kho lưu trữ, bộ mô
phỏng dữ liệu và pipeline học máy.

\begin{table}[H]
\centering
\caption{Nhóm đặc trưng telemetry của Edge Node}
\label{tab:telemetry-groups}
\begin{tabularx}{\textwidth}{L{3cm}L{2.2cm}Y}
\toprule
\textbf{Nhóm} & \textbf{Thông điệp} & \textbf{Ví dụ trường}\\
\midrule
Tài nguyên node & system & \texttt{cpu\_usage\_pct}, \texttt{cpu\_load\_1m/5m}, \texttt{memory\_usage\_pct}, \texttt{swap\_usage\_pct}\\
Nhiệt và điện & system, gpu & \texttt{temperature\_c}, \texttt{thermal\_throttling}, \texttt{power\_w}, \texttt{gpu\_temperature\_c}\\
Lưu trữ & system & \texttt{disk\_usage\_pct}, \texttt{disk\_free\_gb}, \texttt{inode\_usage\_pct}\\
Mạng & system & \texttt{network\_rx/tx\_mb\_s}, \texttt{packet\_loss\_pct}, \texttt{mqtt\_publish\_latency\_ms}\\
GPU & gpu & \texttt{gpu\_usage\_pct}, \texttt{gpu\_memory\_used\_mb}, \texttt{gpu\_frequency\_mhz}, \texttt{gpu\_power\_w}\\
Runtime và suy luận & runtime & \texttt{inference\_fps}, \texttt{effective\_fps\_ratio}, \texttt{inference\_p95\_ms}, \texttt{inference\_error\_rate}, \texttt{runtime\_queue\_size}\\
Camera và đầu ra model & runtime & \texttt{camera\_fps}, \texttt{camera\_frame\_drop\_pct}, \texttt{camera\_freeze\_duration\_sec}, \texttt{detection\_count}, \texttt{confidence\_mean}, \texttt{class\_entropy}\\
\bottomrule
\end{tabularx}
\end{table}

Telemetry \texttt{system} được thu bằng \texttt{psutil}; telemetry \texttt{gpu} được
thu bằng \texttt{nvidia-smi} trên máy có GPU rời (collector trả về rỗng khi công
cụ không có, ví dụ trên Jetson); telemetry \texttt{runtime} do AI Runtime tự đo.

Nhóm \emph{đầu ra model} có ý nghĩa đặc biệt: \texttt{confidence\_mean},
\texttt{detection\_count} và \texttt{class\_entropy} cho phép nhận biết suy giảm
chất lượng suy luận (camera mờ, bị che, đổi góc, ánh sáng thay đổi, dữ liệu trôi)
ngay cả khi CPU, GPU, FPS và latency vẫn bình thường --- điều mà giám sát hạ tầng
thuần túy không phát hiện được.

\subsection{Metric hạ tầng với Prometheus}

Trên mỗi Edge Node, \texttt{node-exporter} (metric hệ điều hành),
\texttt{cAdvisor} (metric container) và \texttt{dcgm-exporter} (GPU NVIDIA, tùy
chọn) chỉ lắng nghe trên \texttt{127.0.0.1}. Một Prometheus chạy ở
\emph{agent mode} thu các exporter mỗi 15~s, gắn nhãn \texttt{node\_id},
\texttt{cluster\_id} và đẩy về Prometheus trên Cloud bằng \texttt{remote\_write}
có xác thực. Mô hình đẩy (push) hoạt động được khi Edge nằm sau NAT; WAL của
agent đệm mẫu khi mất mạng và gửi bù khi kết nối lại.

Trên Cloud, Prometheus lưu metric (mặc định 30 ngày) và đánh giá luật cảnh báo;
Alertmanager định tuyến cảnh báo; Grafana hiển thị và cho phép nhúng panel vào
Dashboard. Bộ luật cảnh báo Edge được liệt kê trong Bảng~\ref{tab:alert-rules}.

\begin{table}[H]
\centering
\caption{Luật cảnh báo metric hạ tầng Edge}
\label{tab:alert-rules}
\begin{tabularx}{\textwidth}{L{4.2cm}Y C{1.6cm}C{1.8cm}}
\toprule
\textbf{Cảnh báo} & \textbf{Điều kiện} & \textbf{Kéo dài} & \textbf{Mức}\\
\midrule
EdgeNodeDown & Node từng báo cáo trong 1 ngày nhưng không đẩy metric trong 5 phút & 2 phút & critical\\
EdgeDiskAlmostFull & Phân vùng gốc sử dụng trên 90\% & 10 phút & warning\\
EdgeHighMemory & Bộ nhớ sử dụng trên 90\% & 10 phút & warning\\
EdgeContainerRestarting & Container khởi động lại hơn 2 lần trong 15 phút & -- & warning\\
EdgeGpuHot & Nhiệt độ GPU trên 85\,°C & 5 phút & warning\\
\bottomrule
\end{tabularx}
\end{table}

\subsection{Thiết kế mô-đun cảnh báo sớm}

Trên nền dữ liệu telemetry đã chuẩn hóa, đồ án thiết kế mô-đun cảnh báo sớm
suy giảm (early warning) theo luồng ở Hình~\ref{fig:monitoring-incident-flow}.

\begin{figure}[H]
\centering
\resizebox{\textwidth}{!}{%
\begin{tikzpicture}[node distance=0.6cm and 0.6cm]
  \node[flowbox] (t) {Telemetry\\5 s/mẫu};
  \node[flowbox, right=of t] (s) {Lưu trữ\\chuỗi thời gian};
  \node[flowbox, right=of s, fill=orange!8] (w) {Cửa sổ trượt\\5 phút (60 mẫu)};
  \node[flowbox, right=of w] (f) {Đặc trưng thời gian};
  \node[flowbox, right=of f, fill=green!6] (m) {Isolation Forest\\XGBoost};
  \node[flowbox, right=of m, fill=red!6] (a) {Risk score\\Cảnh báo};
  \draw[flowarrow] (t) -- (s);
  \draw[flowarrow] (s) -- (w);
  \draw[flowarrow] (w) -- (f);
  \draw[flowarrow] (f) -- (m);
  \draw[flowarrow] (m) -- (a);
\end{tikzpicture}}
\caption{Luồng xử lý của mô-đun cảnh báo sớm suy giảm}
\label{fig:monitoring-incident-flow}
\end{figure}

\textbf{Đặc trưng thời gian.} Với chu kỳ lấy mẫu 5~s và cửa sổ 5 phút, mỗi chỉ
số cốt lõi được biến đổi thành \texttt{current}, \texttt{mean}, \texttt{min},
\texttt{max}, \texttt{std}, \texttt{slope} (hệ số hồi quy tuyến tính) và
\texttt{delta} (chênh lệch so với đầu cửa sổ). Bộ chỉ số khởi đầu gồm khoảng 20
trường: CPU, bộ nhớ và tốc độ tăng bộ nhớ, GPU, nhiệt độ và tốc độ tăng nhiệt, đĩa,
RTT, mất gói, FPS và \texttt{effective\_fps\_ratio}, p95 latency, tỷ lệ lỗi, kích
thước hàng đợi, số lần restart, FPS camera, tuổi khung, tỷ lệ drop, số detection và
confidence trung bình. Một số đặc trưng dẫn xuất như
\texttt{latency\_ratio} = p95 hiện tại / p95 baseline và tốc độ tăng hàng đợi giúp
mô hình học quan hệ đa biến thay vì ngưỡng cứng.

\textbf{Nhãn trạng thái.} Mỗi thời điểm được gắn một trong bốn trạng thái vận
hành: NORMAL, WARNING, DEGRADED, CRITICAL. Ví dụ DEGRADED khi FPS giảm rõ, latency
và hàng đợi tăng; CRITICAL khi runtime FAILED, camera mất kết nối hoặc node OFFLINE.

\textbf{Hai bài toán tách biệt.} Isolation Forest trả lời ``hành vi hiện tại có bất
thường không'' (không cần nhãn). XGBoost trả lời ``trong 10 phút tới có suy giảm
không'': nhãn \texttt{future\_degradation}=1 nếu node chuyển sang DEGRADED hoặc
CRITICAL trong 10 phút sau cửa sổ quan sát. Để tránh rò rỉ nhãn, các trường chỉ
xuất hiện sau sự cố (ví dụ \texttt{runtime\_status}=FAILED) không được dùng làm
đầu vào dự báo. Dữ liệu thô luôn được giữ lại để có thể thay đổi cách trích đặc
trưng mà không cần thu thập lại.

Ví dụ phân biệt: khi GPU tăng, nhiệt độ tăng, latency và hàng đợi tăng trong khi
FPS và confidence giảm, mô hình cảnh báo nguy cơ suy giảm dù từng chỉ số chưa
vượt ngưỡng. Ngược lại, khi GPU tăng do cảnh có nhiều đối tượng nhưng latency,
FPS, hàng đợi và confidence ổn định, trạng thái vẫn là NORMAL, tránh cảnh báo sai.
Trong phạm vi đồ án, hợp đồng telemetry, thu thập và lưu trữ đã được hiện thực;
việc huấn luyện và tích hợp mô hình cảnh báo sớm vào luồng cảnh báo được trình
bày ở mức thiết kế và là hướng phát triển tiếp theo.

\section{Thiết kế AI Assistant}

AI Assistant là một dịch vụ độc lập giúp người vận hành truy vấn và điều khiển
đội thiết bị bằng ngôn ngữ tự nhiên. Nguyên tắc thiết kế cốt lõi là \textbf{Assistant
không bao giờ truy cập trực tiếp cơ sở dữ liệu hay broker}: mọi dữ liệu và thao
tác đều đi qua MCP server của Cloud, nơi cùng dùng các application service với
REST API. Nhờ đó quy tắc nghiệp vụ, kiểm tra dữ liệu và xác thực chỉ được hiện
thực một lần.

\subsection{Công cụ qua Model Context Protocol}

Cloud công bố 17 công cụ qua Model Context Protocol \cite{mcp} (streamable HTTP tại \texttt{/mcp}, bảo vệ bằng
API key), chia làm ba nhóm:

\begin{itemize}
  \item \textbf{Fleet:} \texttt{list\_clusters}, \texttt{get\_cluster},
        \texttt{create\_cluster}, \texttt{list\_nodes}, \texttt{get\_node},
        \texttt{register\_node}, \texttt{list\_runtimes}, \texttt{get\_runtime};
  \item \textbf{Deployment:} \texttt{list\_deployments}, \texttt{get\_deployment},
        \texttt{create\_deployment}, \texttt{update\_deployment},
        \texttt{remove\_deployment}, \texttt{control\_runtime};
  \item \textbf{Monitoring:} \texttt{get\_latest\_telemetry},
        \texttt{get\_telemetry\_history}, \texttt{get\_node\_events}.
\end{itemize}

Schema tham số của công cụ được sinh tự động từ chữ ký hàm Python, và phía
Assistant chuyển mỗi công cụ MCP thành một \texttt{StructuredTool} của LangChain.
Công cụ chỉ được định nghĩa một lần tại Cloud; thêm công cụ mới không cần sửa
Assistant.

\subsection{Kiểm soát thao tác thay đổi}

Các công cụ có tiền tố \texttt{create\_}, \texttt{update\_}, \texttt{remove\_},
\texttt{register\_}, \texttt{control\_} làm thay đổi trạng thái hệ thống. Chúng bị
ẩn hoàn toàn khỏi mô hình trừ khi cấu hình cho phép ghi. Khi được phép, prompt hệ
thống và mô tả của từng công cụ yêu cầu mô hình chỉ gọi sau khi đã mô tả chính
xác hành động và người dùng xác nhận rõ ràng ở một tin nhắn trước đó. Như vậy
có hai lớp bảo vệ: lớp cấu hình (quyết định công cụ có tồn tại với mô hình hay
không) và lớp hội thoại (bắt buộc xác nhận).

\subsection{Đồ thị tác tử}

Một lượt hội thoại được xử lý bằng đồ thị LangGraph \cite{langgraph} bốn nút,
trong đó nút đầu tiên theo mẫu ReAct \cite{yao2023react}
(Hình~\ref{fig:assistant-graph}).

\begin{figure}[H]
\centering
\resizebox{\textwidth}{!}{%
\begin{tikzpicture}[node distance=0.6cm and 0.9cm]
  \node[flowbox, fill=green!6] (start) {Câu hỏi};
  \node[flowbox, right=of start] (react) {\textbf{react}\\gọi tool (ReAct)};
  \node[flowbox, right=of react] (resp) {\textbf{response}\\viết câu trả lời};
  \node[flowbox, right=of resp] (chart) {\textbf{genchart}\\sinh biểu đồ};
  \node[flowbox, right=of chart] (sugg) {\textbf{suggestion}\\gợi ý tiếp};
  \node[flowbox, below=of react, fill=gray!10] (end) {Kết thúc\\(trả lời thẳng)};
  \node[flowbox, above=of react, fill=orange!8] (mcp) {MCP tools\\(Cloud)};
  \draw[flowarrow] (start) -- (react);
  \draw[flowarrow] (react) -- (resp);
  \draw[flowarrow] (resp) -- (chart);
  \draw[flowarrow] (chart) -- (sugg);
  \draw[flowarrow] (react) -- node[right, font=\scriptsize]{không cần tool} (end);
  \draw[flowarrow, <->] (react) -- node[right, font=\scriptsize]{tối đa 6 vòng} (mcp);
\end{tikzpicture}}
\caption{Đồ thị xử lý một lượt hội thoại của AI Assistant}
\label{fig:assistant-graph}
\end{figure}

\begin{itemize}
  \item \textbf{react:} mô hình chính (mặc định Claude Sonnet) lặp tối đa 6 vòng
        \emph{gọi tool $\rightarrow$ nhận kết quả}, chỉ thu thập dữ liệu mà không
        viết câu trả lời. Nếu ngay vòng đầu mô hình không cần tool (chào hỏi, hỏi
        lại khi thiếu đối tượng, hỏi xác nhận), câu trả lời đó được dùng luôn và đồ
        thị kết thúc.
  \item \textbf{response:} tổng hợp kết quả tool thành câu trả lời, phát dần từng
        token về giao diện.
  \item \textbf{genchart:} mô hình nhỏ (Claude Haiku) quyết định dữ liệu có đáng
        vẽ biểu đồ không và sinh đặc tả biểu đồ có cấu trúc; lỗi ở bước này chỉ làm
        mất biểu đồ, không làm hỏng lượt hội thoại.
  \item \textbf{suggestion:} sinh các câu hỏi gợi ý tiếp theo.
\end{itemize}

\textbf{Thực thi tool an toàn.} \texttt{ToolRunner} chạy các lời gọi tool của cùng
một vòng song song (tối đa 8), gộp lời gọi trùng tên và tham số, áp thời gian chờ
60~s cho mỗi tool. Tool không tồn tại, lỗi hoặc quá thời gian được trả về mô hình
dưới dạng dữ liệu lỗi để mô hình tự sửa tham số hoặc giải thích, thay vì làm
dừng cả lượt. Lịch sử hội thoại được giữ theo \texttt{session\_id} bằng checkpointer
của LangGraph, cho phép câu hỏi nối tiếp như ``còn node 2 thì sao?''.

\textbf{Giao thức phát trực tuyến.} Client gửi \texttt{POST /chat/stream} để khởi
tạo lượt và nhận \texttt{stream\_id}, sau đó mở \texttt{GET /chat/stream/\{id\}}
dạng Server-Sent Events. Các sự kiện theo thứ tự: \texttt{status} (đang phân tích,
đang gọi tool nào), \texttt{tool\_end}, \texttt{token}, \texttt{done},
\texttt{chart}, \texttt{suggestions}, hoặc một sự kiện \texttt{error} duy nhất.

\subsection{So sánh với hướng tiếp cận RAG}

Trong bài toán vận hành, câu hỏi chủ yếu hướng tới \emph{trạng thái hiện tại}
của hệ thống (node nào offline, runtime nào lỗi, FPS vừa giảm bao nhiêu). Dữ
liệu này thay đổi liên tục nên truy xuất qua tool có cấu trúc chính xác hơn so
với tìm kiếm văn bản đã lập chỉ mục \cite{lewis2020rag}. Vì vậy đồ án ưu tiên
tool calling qua MCP; RAG trên tài liệu, runbook được giữ ở vai trò mở rộng cho
các câu hỏi về quy trình và kiến thức ổn định.

\section{Kết chương}

Chương 4 đã chuyển yêu cầu của đề tài thành kiến trúc Cloud Control Plane --
Edge Data Plane. Cloud lưu desired state dưới dạng Deployment và phát lệnh qua
MQTT; Edge Agent hội tụ trạng thái container về desired state bằng các thao tác
idempotent, vòng reconcile và desired state lưu trên đĩa; AI Runtime chạy pipeline
camera--suy luận với hàng đợi giới hạn và tự đo chỉ số. Hợp đồng dùng chung,
cơ chế \texttt{last\_command\_id}, chống lệnh trùng và giới hạn số lần thử giữ cho
hệ thống nhất quán qua kết nối không tin cậy. Telemetry được đặc tả thống nhất
làm nền cho giám sát và cảnh báo sớm; AI Assistant truy cập hệ thống duy nhất qua
công cụ MCP với cơ chế kiểm soát thao tác ghi. Chương 5 trình bày việc hiện thực
và triển khai các thành phần này.

